\documentclass[11pt]{article}
\usepackage{fontspec}
\usepackage[french]{babel}
\usepackage{amsmath}
\usepackage{amssymb}
\usepackage{amsthm}
\usepackage{mathrsfs}
\usepackage{tikz-cd}
\usepackage{enumitem}

\newcommand{\printedpage}[1]{\marginpar{\footnotesize p.~#1}}
\newcommand{\msp}[1]{\marginpar{\footnotesize [#1]}}
\newcommand{\bD}{\boldsymbol{\Delta}}
\newcommand{\iddots}{\mathinner{\mkern1mu\raise1pt\hbox{.}\mkern2mu\raise4pt\hbox{.}\mkern2mu\raise7pt\hbox{.}\mkern1mu}}

\begin{document}

\printedpage{4}
% head: A. Grothendieck : Dérivateurs, Ch. VII, Édition des Universités de Montpellier II et Paris VII
% foot: Version 22 décembre 2001

\[
\begin{tikzcd}[column sep=small, row sep=small]
 & & & i \arrow[dl] \arrow[dr] & & & \\
 & & i-1 & & i+1 & & \\
 & \iddots \arrow[dl] & & & & \ddots \arrow[dr] & \\
 i-r & & & & & & i+r, % [sic] i+r -> i+s, s est la longueur de la chaîne contenant u_{i+1}
\end{tikzcd}
\]

où $r$ est la longueur de la chaîne structurale contenant $\underline{u}_i$, $s$ celle de la chaîne structurale contenant $\underline{u}_{i+1}$. Cette catégorie est une somme amalgamée $\bD_r \amalg \bD_s$ (recollées suivant leur origine), mais n'est pas isomorphe à une $\bD_j$, p.~ex. faute d'avoir un objet final.

Cette observation prouve a). Pour prouver b), il suffit de prouver que l'ordre total sur $I$ est déterminé de façon intrinsèque, par le choix de \textit{l'origine} $s_0$ de $I$ (parmi ces % [sic] ces -> ses, possessif
deux extrémités). Supposons que l'on ait déjà déterminé intrinsèquement les sommets $s_0$, $s_1$, $\ldots$, $s_i$ ($i \geq 0$). Comment déterminer $s_{i+1}$? (Quand $i < n$.) Soit $I'$ la sous-catégorie pleine de $I$
\msp{page 7}
formée des $j \in I$ qui ne sont pas parmi les $s_j$ \texttt{[plutôt} $s_k$\texttt{]} déjà construits, sauf le dernier $s_i$ (donc $s_i \in I'$, $\operatorname{card} I' \geq 2$). Alors $I'$ est isomorphe à un intervalle standard, avec $s_i$ correspondant à l'origine. Une des catégories $i\backslash I$, $I/i$ est ponctuelle, l'autre isomorphe à une catégorie $\bD_j$, et de plus, $i$ est objet initial ou final de $\bD_j$. Si c'est l'objet initial $0$, on prend pour $s_{i+1}$ le successeur immédiat de $i$ dans $\bD_j$. Si c'est l'objet final $j$, on prend pour $s_{i+1}$ le prédécesseur immédiat $j-1$.

Si on était parti du but $n$ de $I$, au lieu de partir de la source, on aurait trouvé la relation d'ordre total opposée. Affectant l'élément but du numéro d'ordre $0$, on trouve une renumération de $I$, $i$ devenant $n-i$, et les flèches de transition qui étaient directes pour la relation d'ordre initiale, deviennent rétrogrades dans la relation d'ordre opposée, et inversement. On voit ainsi :

\textbf{Proposition 2.} \textit{Si $I$ est un intervalle standard, il existe un unique couple $(I',\theta)$, où $I'$ est un intervalle standard, et $\theta$ un isomorphisme $I \xrightarrow{\sim} I'$ qui transforme le but de $I$ dans la source de $I'$ (et la source de $I$ en le but de $I'$).}

Cet intervalle $I'$ est appelé l'intervalle \textit{opposé} de $I$. Il faut le voir, intrinsèquement, comme la `même' categorie $I$, % [sic] categorie -> catégorie, accent manquant
mais avec une `orientation' opposée.

\msp{page 8}

Une façon plus claire de voir ceci est la suivante. Soit $J$ une catégorie intervalle, i.e. isomorphe à un intervalle standard. Alors il y a exactement (si $I \not\simeq e$) \textit{deux} couples $(I,\theta)$ d'un intervalle standard $I$ et d'un isomorphisme $\theta : I \xrightarrow{\sim} J$. Ces deux isomorphismes induisent sur $J$ deux ordres totaux, opposés l'un de l'autre. Chaque couple est déterminé de façon unique par ledit ordre total, qu'on appelera % [sic] appelera -> appellera, orthographe
une \textit{orientation} de l'intervalle $J$. Ainsi $J$ a deux orientations, et chacune détermine un \textit{isomorphisme orienté} (i.e. respectant
\printedpage{5}
% head: A. Grothendieck : Dérivateurs, Ch. VII, Édition des Universités de Montpellier II et Paris VII
% foot: Version 22 décembre 2001
les orientations) de $J$ avec un intervalle standard (\footnote[2]{\textbf{NB} On parlera de l'intervalle \textit{orienté opposé} $I^{\mathrm{opp}}$ d'un intervalle \textit{orienté} $I$, qu'il faut se garder de confondre avec la catégorie opposée $I^o$. Ici $I^{\mathrm{opp}}$ a même catégorie sous-jacente que $I$, mais avec une orientation opposée.}). Ces deux intervalles sont les mêmes si et seulement s'il existe un automorphisme de $J$ (nécessairement unique) qui échange ses deux orientations, ou ce qui revient au même, qui échange entre elles les deux extrémités. (Ainsi les catégories opposées l'une de l'autre
\[
\begin{tikzcd}[column sep=small, row sep=small]
 & \bullet \arrow[dl] \arrow[dr] & & & & \bullet & \\
 \bullet & & \bullet & & \bullet \arrow[ur] & & \bullet \arrow[ul]
\end{tikzcd}
\]
sont chacunes % [sic] chacunes -> chacune, orthographe
des intervalles symétriques.) Choisir une orientation revient aussi `a % [sic] `a -> à, accent imprimé comme guillemet
choisir l'origine parmi les deux extrémités.

Notons que la \textit{catégorie opposée} à un intervalle est un intervalle, et la catégorie opposée à un intervalle standard est l'intervalle standard qui s'en déduit en changeant tous les signes sur les intervalles de transition $[i-1,i]$ ($1 \leq i \leq n$). On se gardera de confondre la propriété $I \simeq I^o$, avec la propriété de symétrie.

\msp{page 9}

Ainsi les deux catégories ($\Psi_1$ et $\Phi_1$) ci-dessus sont symétriques, mais opposées l'une de l'autre sans être isomorphes. D'autre part, les $\bD_i$ ne sont pas symétriques, mais chacune isomorphe à la catégorie opposée.

Si $I$ est isomorphe à la catégorie opposée $I^o$, et si $\theta$ est un tel isomorphisme (nécessairement unique), alors il est immédiat que $\theta$ renverse l'orientation (sinon ce serait l'identité sur les sommets, et les flèches de transition seraient à la fois directes et indirectes), donc $\theta(0) = n$, $\theta(1) = n-1$, et on veut que $\underline{u}_1$ et $\underline{u}_n$ sont de même signe. Mais si $I$ est symétrique, $\underline{u}_1$ et $\underline{u}_n$ sont de signe opposé. Donc $I$ ne peut être à la fois symétrique et isomorphe à l'intervalle opposé \texttt{[plutôt à la catégorie opposée; exception :} $I = \{0\}$\texttt{]}.

Pour toute catégorie $\bD_i$, celle-ci étant un intervalle, il y a (si $i > 0$) \textit{deux orientations.} En tant qu'intervalles orientés, il y a $\bD_i^+$ (avec l'orientation ordinaire, donné par l'ordre ordinaire sur $\operatorname{Ob} \bD_i = [0,i]$), et $\bD_i^-$ avec l'orientation opposée.

\printedpage{6}
% head: A. Grothendieck : Dérivateurs, Ch. VII, Édition des Universités de Montpellier II et Paris VII
% foot: Version 22 décembre 2001

\msp{page 10}

\section*{2\quad Sous-catégories pleines d'un intervalle}

Soit $I$ une catégorie intervalle orientée, de longueur $n \geq 1$. Nous identifierons $\operatorname{Ob} I$ à l'ensemble des entiers $[0,n]$. Pour une chaîne $\underline{u}_i$, $\underline{u}_{i+1}$, $\ldots$, $\underline{u}_{i+l-1}$ de flèches de transition dans $I$, on appelle \textit{domaine} de la chaîne, l'ensemble des sources et de buts de ces éléments. C'est donc un intervalle d'entiers $[i-1, i+l-1]$, de longueur $l$ égale à celle de la chaîne, dont l'origine est déterminé % [sic] déterminé -> déterminée, accord avec l'origine
par la première flèche $\underline{u}_i$ de la chaîne (c'est sa source, si $\underline{u}_i$ est directe, et son but si $\underline{u}_i$ est rétrograde), et dont l'extrémité est déterminée par la dernière $\underline{u}_{i+l-1}$. Les intervalles d'entiers (vus comme des sommets de $I$) correspondant aux chaînes structurales sont appelés les \textit{intervalles structuraux} ou \textit{intervalles critiques} de $I$. Contrairement aux chaînes structurales, ces intervalles ne forment pas une partition de $I$, puisque deux intervalles structuraux peuvent avoir une extrémité commune. Si les chaînes structurales sont
\[
C_1,\ C_2,\ \ldots\ ,\ C_r \qquad (r = \text{ longueur utile de } I)
\]
de longueurs
\[
l_1,\ l_2,\ \ldots\ ,\ l_r,
\]
alors les intervalles structuraux sont
\[
\underbrace{[0,l_1]}_{I_1},\ \underbrace{[l_1,l_1+l_2]}_{I_2},\ \ldots\ ,\ \underbrace{[l_1+\cdots+l_{r-1},l_1+\cdots+l_{r-1}+l_r]}_{I_r},
\]
une suite d'intervalles d'entiers dont deux consécutifs ont une extrémité commune, et tels que si $i < j$, chaque élément de $I$ % [sic] I -> I_i, comparaison des éléments de I_i et de I_j
soit $\leq$ chaque élément de $I_j$. On appelle \textit{sommet critique} de $I$ les extrémités des intervalles structuraux.

\textbf{Proposition 3.} \textit{Soit $s \in I$, soit $I' = I/s$ (sous-catégorie pleine de $I$ des $s'$ tels qu'il existe $s' \longrightarrow s$), $I'' = I/s$ \textup{\texttt{[plutôt} $I'' = s\backslash I$\texttt{]}}, $I'_0 = (I/s)\backslash\{s\}$, $I''_0 = (s\backslash I)\backslash\{s\}$} (\footnote[3]{i.e. qu'on ait : il n'y a pas de flèches entre un élément de $I'_0$ et un élément de $I''_0$.}). Pour que $s$ soit critique, il faut et il suffit que $I\backslash\{s\}$ soit catégorie somme de $I'_0$ et $I''_0$.

\msp{page 11}

\textbf{Proposition 4.} \textit{Soient $i,j$ dans $I$, avec $i < j$. Conditions équivalentes :}
\begin{enumerate}
\item[\textit{a)}] \textit{Il existe une flèche entre $i$ et $j$ (soit $i \longrightarrow j$, soit $j \longrightarrow i$).}
\item[\textit{b)}] \textit{$i$ et $j$ appartiennent à un même intervalle structural (nécessairement unique).}
\item[\textit{c)}] \textit{La sous-catégorie pleine de $I$ définie par l'ensemble d'objets $[i,j]$ est isomorphe à une $\bD_\alpha$ (on aura $\alpha = j-i$).}
\end{enumerate}

\printedpage{7}
% head: A. Grothendieck : Dérivateurs, Ch. VII, Édition des Universités de Montpellier II et Paris VII
% foot: Version 22 décembre 2001
On dira alors que $i$ et $j$ sont \textit{directement connectés} dans $I$. (Et on utilisera la même terminologie dans une catégorie quelconque : $i$ et $j$ sont directement connectés si et seulement si $\operatorname{Hom}_I(i,j)$ \textit{ou} $\operatorname{Hom}_I(j,i)$ est non vide.)

\textbf{Proposition 5.} \textit{Soit $J \subseteq I$ une partie de $I$, considérons la sous-catégorie pleine de $I$ qu'elle définit. Alors :}
\begin{enumerate}
\item[\textit{a)}] \textit{Les composantes connexes $J_\alpha$ de $J$ sont des intervalles et l'ordre induit sur un $J_\alpha$ par celui de $I$ est une orientation (dite \textup{orientation induite}) de $J_\alpha$.}
\item[\textit{b)}] \textit{Supposons $\operatorname{card} J > 1$. Soit $i$ le plus petit élément de $J$, $j$ le plus grand élément, donc $i < j$. Pour que $J$ lui-même soit un intervalle, il faut et il suffit que tout sommet critique dans $[i,j]$ soit élément de $J$.}
\end{enumerate}

(\footnote[4]{On appelle \textit{sous-intervalle} de $I$ une sous-categorie % [sic] sous-categorie -> sous-catégorie, accent manquant
pleine connexe, i.e. qui est elle-même un intervalle. Notons d'ailleurs que \textit{toute} sous-catégorie $I'$ de $I$ qui est un intervalle, est nécessairement pleine, cf. plus bas \texttt{[paragraphe]} des foncteurs entre intervalles.}).

\textbf{NB} Si $s \in [i,j]$ est un sommet critique de $I$ qui n'est pas dans $J$, alors par la proposition 3 il disconnecterait $J$, i.e. $J$ serait somme directe de deux sous-catégories $J'$ et $J''$ non vides ($i \in J'$, $j \in J''$), donc $J$ ne serait pas un intervalle.

\msp{page 12}

\textbf{Corollaire.} \textit{L'ensemble des sommets critiques de $I$ est aussi l'ensemble des sommets du plus petit sous-intervalle de $I$ contenant les extrémités. Plus généralement, si $i < j$ sont deux sommets de $I$, le plus petit sous-intervalle de $I$ contenant $i$, $j$ est formé de $i$, $j$ et des sommets critiques $s$ tels que $i < s < j$.}

\textbf{Définition.} Un intervalle est dit \textit{réduit} (ou \textit{irrédondant}) si tous ses sommets sont critiques, ou ce qui revient au même, si ses chaînes structurales sont de longueur 1.

Un intervalle orienté \textit{réduit} est connu, à isomorphisme près, quand on connaît
\begin{enumerate}
\item[a)] sa longueur $\mathrm{n}$,
\item[b)] sa direction initiale $+$ (directe) ou $-$ (rétrograde).
\end{enumerate}

\textbf{NB} Cette direction est la même que la direction finale si $n$ est impair, elle est opposée si $n$ est pair.

Ainsi, suivant les deux cas de direction initiale, on trouve
\[
\begin{aligned}
0 \longrightarrow 1 \longleftarrow 2 \longrightarrow \cdots \qquad\qquad \cdots\ n & \qquad (\text{Avec } & n-1 \longrightarrow n & \quad \text{si } n \text{ impair,}\\
& & n-1 \longleftarrow n & \quad \text{si } n \text{ pair.)}\\
0 \longleftarrow 1 \longrightarrow 2 \longleftarrow \cdots \qquad\qquad \cdots\ n & \qquad (\text{Avec } & n-1 \longleftarrow n & \quad \text{si } n \text{ impair,}\\
& & n-1 \longrightarrow n & \quad \text{si } n \text{ pair.)}
\end{aligned}
\]

\printedpage{8}
% head: A. Grothendieck : Dérivateurs, Ch. VII, Édition des Universités de Montpellier II et Paris VII
% foot: Version 22 décembre 2001

\textbf{Proposition 6.} \textit{Un sous-intervalle $J$, ayant $i$ et $j$ ($i < j$) comme plus petit et plus grand sommet, est réduit si et seulement si tout $s \in J$ distinct de $i$ et de $j$ est critique.}

\textbf{Corollaire.} \textit{La sous-catégorie pleine de $I$ formée des sommets critiques est un sous-intervalle. On l'appelle l'intervalle réduit associé à $I$, noté $I_{\textup{réd}}$.}

\end{document}
